\section{Experiments}
\label{sec:experiments}
\subsection{Experimental setup}
\label{subsec:experimental_setup}

All experiments were performed using Cooja with the same simulation configuration.
The radio medium was configured as \gls{udgm}, each mote was started with a delay of 1 second, and the random seed was fixed to \texttt{123456} in order to ensure reproducibility.

The experiments were based on the \gls{rpl}-\gls{udp} example included in Contiki-NG.
The clients and server used in the simulations are Cooja motes implemented in \texttt{examples/rpl-udp}, with the server acting as the \gls{dodag} root and the clients participating as \gls{rpl} nodes.
Attack-specific motes were added when required by the scenario.

For each experiment, we collected the Cooja \texttt{.csc} configuration, a screenshot of the simulated network, the \gls{dag} reconstructed from the execution logs, the complete mote log, and the two \gls{csv} files generated by the Python \glspl{ids}.

\subsection{Experimental scenarios}
\label{subsec:experimental_scenarios}

The experimental campaign consists of one baseline scenario and eight attack scenarios. In the following table, \textit{B} denotes a Blackhole node and \textit{S} denotes a Sinkhole (client) node.
The position refers to the attacker's position within the resulting network topology.

\begin{table}[h]
\centering
\caption{Experimental scenarios}
\label{tab:experimental_scenarios}
\begin{tabular}{lcccc}
\hline
\textbf{Scenario} & \textbf{Server} & \textbf{Clients} &
\textbf{Attackers} & \textbf{Position / Configuration} \\
\hline
Baseline & 1 & 11 & -- & -- \\
Blackhole 1 & 1 & 10 & 1 B & Middle \\
Blackhole 2 & 1 & 10 & 1 B & Leaf \\
Blackhole 3 & 1 & 10 & 1 B & Bottleneck \\
Sinkhole 1 & 1 & 10 & 1 S & Middle \\
Sinkhole 2 & 1 & 10 & 1 S & Leaf \\
Sinkhole 3 & 1 & 10 & 1 S & Bottleneck \\
Sinkhole 4 & 1 & 10 & 1 S & Original MRHOF + hysteresis \\
Blackhole 4 & 1 & 10 & 1 B+S & Combined with Sinkhole \\
\hline
\end{tabular}
\end{table}
